% SOURCES: docs/0.2.0v/plan.md (§1 North Star, §4 Guiding Principles);
%          README.md (Key invariants); CLAUDE.md (Core Design Rules)
% GUIDANCE: One clear primary objective + a short list of specific, measurable
% sub-objectives. Keep each sub-objective testable so ch. 10–11 can report on it.

\chapter{Project Objective}
\label{ch:objective}

\section{Primary Objective}
% DRAFT (review):
To design and build a desktop application that lets non-developers author modern,
semantic, responsive, and accessible web pages by direct manipulation on a canvas,
and that exports a portable, human-readable HTML/CSS bundle with opt-in,
independently toggleable JavaScript behaviours.

\section{Specific Objectives}
% DRAFT (review): these map directly onto measurable outcomes in ch. 10–11.
\begin{enumerate}[label=\textbf{O\arabic*.}, leftmargin=3.2em]
  \item \textbf{Single source of truth.} Represent every page as one Document
        Model; the canvas renders it and the generator walks it, with no
        component owning state independently.
  \item \textbf{Semantic, deterministic output.} Emit semantic HTML5 by role and
        guarantee byte-identical output for identical input (no timestamps, no
        random identifiers).
  \item \textbf{Accessibility by construction.} Enforce a hard gate that blocks
        export on any critical or serious accessibility violation; honour a single
        \texttt{<h1>}, heading order, \texttt{alt} text, ARIA labels, focus
        visibility, and reduced-motion preferences.
  \item \textbf{Tokens-driven, responsive styling.} Drive all styling through
        design tokens with light/dark theming and per-breakpoint values
        (desktop/tablet/mobile/small).
  \item \textbf{Opt-in, vetted runtime.} Ship only the JavaScript the author
        enables; with all behaviours disabled, produce JavaScript-free output.
  \item \textbf{Accessible authoring surfaces (v1.0.0).} Provide draw-to-create,
        match-layout, and a headless MCP server that reuse the same operations and
        export pipeline.
\end{enumerate}

\section{Scope}
% SOURCES: docs/0.2.0v/plan.md (§2 Scope Tiers, §22 Out of Scope); draft/Features/*.md
% DRAFT (review):
The project scope is the \emph{Tier~1} output-capability set (see
\texttt{draft/Features/features.md}). Editor-experience niceties (Tier~2) and
cloud/accounts/collaboration (Tier~3) are out of scope, except for the three
Tier~2 authoring surfaces consciously brought into v1.0.0 (draw, match, MCP).
% TODO: state the boundary explicitly for your defence.
